\section{相关理论基础}
\subsection{类别不均衡与类别偏置}
类别不均衡指数据集中不同类别的样本数量差距较大，分为天然不均衡、人为采样不均衡等情况。类别不均衡往往会带来类别偏置问题：模型在训练时会偏向样本数量多的大类，在少数类别上决策边界发生偏移，预测输出存在系统性偏差。

类别偏置和单纯的数据不平衡有区别：即使样本数量均衡，若各类特征分布复杂度不同，依然会产生类别偏置。常用评价指标包括召回率、F1值、均衡准确率、偏差方差分解等。在不均衡分类任务中，不能只依靠总体准确率评价模型效果，需要重点关注各类别上的性能差异。

\subsection{神经网络特征传播机制}
深度神经网络由多层非线性变换堆叠而成，输入样本经过逐层卷积、激活、池化或者全连接运算，特征不断被提取、映射、抽象。浅层提取边缘、纹理等底层特征，深层学习语义、类别相关的高层特征。

在训练迭代过程中，梯度反向传播不断更新网络参数，各类特征的分布持续变化。类别偏置也会随着迭代逐步累积：大类样本的梯度主导参数更新，不断挤压少数类的特征空间，最终造成偏置不断放大。

\subsection{波动力学基础}
波动力学以波动方程为核心，描述物理量随时间、空间的传播与演化规律。一维波动方程形式为：
\[
\frac{\partial^2 u}{\partial t^2}=c^2\frac{\partial^2 u}{\partial x^2}
\]
其中 $u(x,t)$ 为波动状态量，$c$ 为波速。波动具有传播、叠加、衰减、色散等特性，状态会随时间逐步扩散、演化。

本研究借鉴波动演化思想，不直接求解物理波动方程，而是将类别偏差视作一种“波动信号”：把训练迭代当作时间维度，类别之间的特征相似度当作空间连接，刻画偏差在训练过程中传播、累积、衰减的规律，以此实现对类别偏置变化趋势的预测。

\subsection{不均衡学习常用方法}
1. 数据层面：重采样，包括对少数类过采样、对多数类欠采样；生成式方法利用GAN、扩散模型合成少数样本。
2. 损失层面：加权交叉熵、Focal Loss、LDAM等，通过调整各类损失权重抑制偏置。
3. 训练策略层面：解耦训练、两阶段训练、迁移学习，先学习通用特征再微调少数类。

上述方法大多是补偿策略，缺少对偏置演化趋势的预测能力。本文正是基于这一点，引入波动演化模型提前预测偏置，为动态补偿提供依据。
